\begin{proposition}[6.7.4]
\label{IV.2.s6.7.4-fr}
Para un preesquema localmente noetheriano $Z$ y un punto $z\in Z$,
sea $\boldsymbol{P}(Z,z)$ una de las propiedades siguientes:
\begin{enumerate}
\item[\textup{a)}] $Z$ es regular en el punto $z$.
\item[\textup{b)}] $Z$ verifica $(R_n)$ en el punto $z$.
\item[\textup{c)}] $Z$ es normal en el punto $z$.
\item[\textup{d)}] $Z$ es reducido en el punto $z$.
\end{enumerate}
Dicho esto, bajo las hipótesis y con las notaciones de \textup{(6.7.1)}, si
$\boldsymbol{P}(X',x')$ es verdadera, lo mismo ocurre con $\boldsymbol{P}(X,x)$. La recíproca es
verdadera si $k'$ es una extensión separable de $k$.
\end{proposition}

\begin{proof}
El caso en que $\boldsymbol{P}$ es la propiedad \emph{d)} solo se ha incluido a título de
recordatorio, pues ya ha sido tratado en (2.1.13 y 4.6.1). El mismo
razonamiento que al comienzo de la demostración de (6.7.1) muestra que la
primera afirmación resulta, respectivamente, de (6.5.2, (i)), (6.5.3, (i))
y (6.5.4, (i)); del mismo modo, la segunda afirmación resultará de (6.5.2,
(ii)), (6.5.3, (ii)) y (6.5.4, (ii))

\oldpage[IV-2]{147}

con tal de que se demuestre que $p^{-1}(x)$ es un \emph{preesquema regular}; dicho
de otro modo, la cuestión se reduce a demostrar el
\end{proof}

\begin{lemma}[6.7.4.1]
\label{IV.2.s6.7.4.1-fr}
Sean $K$, $L$ dos extensiones de un cuerpo $k$, una de las cuales es de tipo
finito. Si $L$ es una extensión separable de $k$, el anillo $K\otimes_kL$
es regular.
\end{lemma}

\begin{proof}
Distingamos dos casos:

\begin{enumerate}
\item[A)] $L$ es una extensión \emph{de tipo finito} de $k$. Entonces, con las
notaciones de (6.7.1.1), se puede suponer que $L$ es una extensión \emph{finita
separable} de $k(\mathbf t)$. Para todo $s\in\operatorname{Spec}(A)$,
$\boldsymbol{k}(s)\otimes_{k(\mathbf t)}L$ es entonces producto directo de un
número finito de cuerpos, extensiones finitas separables de $\boldsymbol{k}(s)$, y, por
tanto, un anillo \emph{regular}; resulta entonces de (6.5.2, (ii)) que basta
demostrar que el anillo $A$ es regular; como $A=S^{-1}A'$, basta demostrar
que $A'$ es regular (\textbf 0, 17.3.6), pero ya se vio que así era en la demostración
de (6.7.1.1).
\item[B)] Sea $(L_\lambda)$ la familia filtrante de las subextensiones de
$L$ que son de tipo finito sobre $k$; en virtud de A), cada uno de los
anillos $C_\lambda=K\otimes_kL_\lambda$ es regular; por otra parte, para
$L_\lambda\subset L_\mu$, $L_\mu$ es un $L_\lambda$-módulo plano, por lo
que $C_\mu$ es un $C_\lambda$-módulo plano $(\mathbf{0}_{\mathrm I},6.2.1)$. Como
$C=K\otimes_kL=\varinjlim(K\otimes_kL_\lambda)$ es noetheriano, se puede
aplicar (5.13.7), y $C$ es, por tanto, regular.
\end{enumerate}
\end{proof}
